% -----------------------------------------------------------------------------
% This file is automatically generated.
% Do not edit manually.
% Source: notebooks/support/regression-analysis/support.Rmd
% -----------------------------------------------------------------------------


\noindent Next, we set the audit parameters and construct the acceptable difference range. The acceptable difference range is expressed around zero, because it is applied to the difference between the recorded amount and the expectation.

\par\addvspace{\topsep}
\begingroup
\fvset{listparameters={%
  \setlength{\topsep}{0pt}%
  \setlength{\partopsep}{0pt}%
  \setlength{\parsep}{0pt}%
  \setlength{\itemsep}{0pt}%
}}
\begin{Verbatim}[breaklines=true,formatcom=\color{ada_blue}]
# Set audit parameters
PM = 15000000         # Performance materiality
beta = 0.01           # Efficiency risk used for the acceptable difference range
df = df_res
# Standardized materiality
delta = PM / annual_se
# Critical values for a two-sided 99% acceptable difference range
c_L = t.ppf(beta / 2, df=df)
c_U = t.ppf(1 - beta / 2, df=df)
# Acceptable difference range for the annual difference
ADR_lower = c_L * annual_se
ADR_upper = c_U * annual_se
cat("Acceptable Difference Range for the annual difference:\n")
cat("Lower bound:", ADR_lower, "\n")
cat("Upper bound:", ADR_upper, "\n")
\end{Verbatim}
\begin{Verbatim}[breaklines=true,formatcom=\color{ada_light_blue}]
Acceptable Difference Range for the annual difference:

Lower bound: -17389050.5

Upper bound: 17389050.5
\end{Verbatim}
\endgroup
\par\addvspace{\topsep}
